% ID: SEL-D0017-BD-E08
% ESTADO: cotejada. Segunda pasada visual completa y borrador compilado revisados.
% Fuente: Selectividad/matematicas_aplicadas/matematicas_aplicadas_P_E.pdf
% SHA-256: bc88bfeb0f184eb0c806c4efe04e84c24f3e5fcef9caaa12e94fd48d9f67a389
% Página: 2; bloque D; ejercicio original 8.

Una tienda de ropa quiere estudiar la aceptación de un nuevo sistema de
pago a través del teléfono móvil. Para ello realiza una encuesta entre 200
de sus clientes elegidos al azar, resultando que 150 de ellos sí estarían
dispuestos a usar el nuevo sistema de pago.
\begin{enumerate}[label=\alph*)]
  \item \textbf{(1.5 puntos)} Determine un intervalo de confianza al 97\%
  para estimar la proporción de clientes de esa tienda que estarían
  dispuestos a usar el nuevo sistema de pago.
  \item \textbf{(1 punto)} Mediante una nueva encuesta se quiere estimar
  la proporción de clientes de esa tienda que usarían el nuevo sistema de
  pago, con un error máximo del 3\% y un nivel de confianza del 94\%.
  Suponiendo que se mantiene la proporción muestral del apartado anterior,
  \textquestiondown{}a cuántos clientes como mínimo habría que realizar
  la encuesta?
\end{enumerate}
